\begin{abstract}

Adaptive evolution is shaped by the interplay between natural selection and the generation of phenotypic variation, which is itself shaped by underlying developmental mechanisms.
Despite the centrality of this relationship, how variation-generating processes interact with complex adaptive landscapes remains poorly understood.
To address this, we embedded a set of mechanistic developmental models, composed of simple regulatory motifs, within adaptive landscapes of varying complexity and examined how developmental bias, i.e. the biased distribution of phenotypic effects of random genetic mutations, and selection jointly determine evolutionary outcomes.
First, we observe that developmental bias not only leads to a transient bias in evolutionary trajectories, but also to a distortion of the domains of attraction surrounding alternative evolutionary optima.
Second, as landscapes become more rugged relative to the scale of phenotypic variation, developmental bias becomes increasingly decisive for which of several optima a population evolves towards.
Thus, developmental mechanisms become increasingly important in predicting evolutionary outcomes the more rugged adaptive landscapes are.
Third, because developmental bias emerges from regulatory dynamics, it can also itself evolve as populations move through phenotype space: populations beginning with nearly identical mutational variation can diverge in their standing and mutational variation, altering the distribution of accessible variation and ultimately reaching different optima.
Altogether, the relative importance of selection and developmental bias is a property of the combination of both, not a fixed feature of either.

% Adaptive evolution is shaped by the interplay between natural selection and the generation of phenotypic variation, which is itself structured by underlying developmental mechanisms.
% Despite the centrality of this relationship, how variation-generating processes interact with complex adaptive landscapes remains poorly understood.
% To address this, we embedded mechanistic developmental models—composed of simple two-gene regulatory motifs—within adaptive landscapes of varying complexity and examined how developmental bias (the anisotropic phenotypic distribution of random genetic mutations) and selection jointly determine evolutionary outcomes.
% First, we show that developmental bias deflects evolutionary trajectories transiently and can permanently distort the domains of attraction surrounding alternative optima.
% Because biased mutational input steers populations along axes of high mutational variance rather than directly uphill along the selection gradient, trajectories are carried across fitness ridges into alternative basins, rotating basin boundaries toward the leading axis of mutational variation ($M_{\max}$).
% Second, as adaptive landscapes become more rugged relative to the scale of phenotypic variation, developmental bias becomes increasingly decisive for which optimum a population reaches: when fitness peaks are smaller and more densely packed, the same bias-induced deflection crosses basin boundaries far more frequently, increasing peak mismatches relative to an unbiased baseline and tightening the alignment between evolutionary divergence and $M_{\max}$.
% Third, because nonlinear gene regulation makes mutational sensitivities state-dependent, the local mutational covariance matrix ($\mathbf{M}$) varies across phenotype space within any given network topology, so a population experiences changing local developmental bias as it moves through its genotype-phenotype map.
% Moreover, different regulatory topologies that generate nearly identical local $\mathbf{M}$-matrices at the same ancestral phenotype can subsequently follow different trajectories in standing ($\mathbf{G}$) and mutational ($\mathbf{M}$) covariation as they traverse their respective maps, reliably settling on different optima.
% Altogether, the relative importance of selection and developmental bias in determining evolutionary outcomes is not a fixed property of either process, but an emergent consequence of how the structure of development interacts with the geometry of the adaptive landscape.

Keywords: \textbf{evolvability, G-matrix, M-matrix, gene regulatory networks, basin of attraction}

\end{abstract}
